\documentclass[10pt,a4paper]{article}
\usepackage[margin=0.6in]{geometry}
\usepackage{amsmath,amssymb,amsfonts,stmaryrd}
\usepackage{xcolor}
\usepackage{enumitem}
\usepackage{tcolorbox}
\usepackage{multicol}
\usepackage{listings}
\usepackage{hyperref}

\definecolor{isabelleblue}{RGB}{0, 70, 140}
\definecolor{isabellekw}{RGB}{120, 0, 120}
\definecolor{boxbg}{RGB}{248, 249, 250}

\tcbset{
  colback=boxbg,
  colframe=isabelleblue,
  fonttitle=\bfseries\large,
  boxrule=0.8pt,
  arc=3pt,
  left=8pt,
  right=8pt,
  top=6pt,
  bottom=6pt
}

\title{\textbf{\huge List Induction (Part 1: Structural Induction Principles)}\\[0.3em]
\large Lecture Notes Transcription}
\author{\textbf{Course:} Interactive Theorem Proving (7CCSACTP)}
\date{}

\begin{document}

\maketitle
\vspace{-1.5em}

\begin{tcolorbox}[title=1. Induction on Lists (Structural Induction)]
\begin{multicols}{2}
\textbf{$\rightarrow$ Context:}
\begin{itemize}[leftmargin=1.2em, itemsep=2pt]
    \item Introduction to recursion on lists
    \item Structural induction principle on finite lists
\end{itemize}

\textbf{$\rightarrow$ Intuition \& Analogy:}
\begin{itemize}[leftmargin=1.2em, itemsep=2pt]
    \item \textbf{Natural numbers:} Built via finite applications of \texttt{Suc}.
    \item \textbf{Lists:} Any list is constructed with a finite sequence of \texttt{Cons} (\texttt{\#}) ending in \texttt{[]}:
    \[ x_1 \# x_2 \# \dots \# x_n \# [] \]
    \item To prove $P(xs)$ holds for all lists $xs$:
    \begin{itemize}[leftmargin=1em, label=$\bullet$]
        \item $\textbf{goal}_1 \textbf{ (Base Case):}$ $P([])$
        \item $\textbf{goal}_2 \textbf{ (Step Case):}$ Assuming $P(xs)$, prove $P(x \# xs)$ for any $x$.
    \end{itemize}
\end{itemize}

\columnbreak

\textbf{Formal Inference Rules:}
\begin{itemize}[leftmargin=1.2em, itemsep=4pt]
    \item \textbf{Natural Numbers:}
    \[ \frac{P(0) \qquad P(m) \Longrightarrow P(m+1)}{P(n)} \]
    
    \item \textbf{Lists:}
    \[ \frac{P([]) \qquad P(ys) \Longrightarrow \forall y.\; P(y \# ys)}{P(xs)} \]
\end{itemize}
\end{multicols}
\end{tcolorbox}

\vspace{0.3em}

\begin{tcolorbox}[title=2. Worked Example: Associativity of List Append]
\textbf{Definition of Append (\texttt{@}):}
\[ \text{append } [] \; ys = ys \qquad \text{append } (x \# xs) \; ys = x \# (\text{append } xs \; ys) \]

\textbf{Theorem:} $(xs \mathbin{@} ys) \mathbin{@} zs = xs \mathbin{@} (ys \mathbin{@} zs) \qquad (\text{Induction on } xs)$

\begin{multicols}{2}
\textbf{Base Case ($xs = []$):}
\begin{itemize}[leftmargin=1.2em, itemsep=2pt]
    \item $\textbf{goal}::$ $([] \mathbin{@} ys) \mathbin{@} zs = [] \mathbin{@} (ys \mathbin{@} zs)$
    \item \textbf{Proof steps:}
    \begin{enumerate}[leftmargin=1.2em, itemsep=1pt]
        \item $([] \mathbin{@} ys) \mathbin{@} zs = ys \mathbin{@} zs$ \quad \textit{(by append.simps)}
        \item $\dots = [] \mathbin{@} (ys \mathbin{@} zs)$ \quad \textit{(by append.simps)}
    \end{enumerate}
    \item \textsf{finally} $([] \mathbin{@} ys) \mathbin{@} zs = [] \mathbin{@} (ys \mathbin{@} zs) \quad \checkmark$
\end{itemize}

\columnbreak

\textbf{Step Case ($xs \to x \# xs$):}
\begin{itemize}[leftmargin=1.2em, itemsep=2pt]
    \item \textbf{Assume I.H.:} $(xs \mathbin{@} ys) \mathbin{@} zs = xs \mathbin{@} (ys \mathbin{@} zs)$
    \item $\textbf{goal}::$ $((x \# xs) \mathbin{@} ys) \mathbin{@} zs = (x \# xs) \mathbin{@} (ys \mathbin{@} zs)$
    \item \textbf{Proof steps:}
    \begin{enumerate}[leftmargin=1.2em, itemsep=1pt]
        \item $((x \# xs) \mathbin{@} ys) \mathbin{@} zs = (x \# (xs \mathbin{@} ys)) \mathbin{@} zs$ \quad \textit{(append.simps)}
        \item $\dots = x \# ((xs \mathbin{@} ys) \mathbin{@} zs)$ \quad \textit{(append.simps)}
        \item $\dots = x \# (xs \mathbin{@} (ys \mathbin{@} zs))$ \quad \textit{(by I.H.)}
        \item $\dots = (x \# xs) \mathbin{@} (ys \mathbin{@} zs)$ \quad \textit{(append.simps)}
    \end{enumerate}
    \item \textsf{finally} $((x \# xs) \mathbin{@} ys) \mathbin{@} zs = (x \# xs) \mathbin{@} (ys \mathbin{@} zs) \quad \checkmark$
\end{itemize}
\end{multicols}
\end{tcolorbox}

\end{document}
